\appendix
\onecolumn % The $\mathtt{\backslash onecolumn}$ command above can be kept in place if you prefer a one-column appendix, or can be removed if you prefer a two-column appendix.  Apart from this possible change, the style (font size, spacing, margins, page numbering, etc.) should be kept the same as the main body.

\section*{Organization of Appendix}

The appendix is organized as follows: 
in Section \ref{sec:fc-deep-dive} further details about failure categories and failure modes are given, 
in Section \ref{appendix:mas-with-annotated-traces-appendix} we provide some details about the multi-agent systems we have annotated and studied, 
in Section \ref{sec:python_library} we describe how to use the \taxonomy{} easily as a python library using \verb|pip install|,
in Section \ref{sec:programdev} we describe the tasks in ProgramDev and ProgramDev-v2 Datsaet,
in Section \ref{appendix:mas-failure-modes-corr} we plot the correlations between MAS failure modes, 
in Section \ref{sec:failure_across_model} we analyze the failure comparison between models and MAS,
in Section \ref{sec:better-mas} we discuss some tactical approaches and structural strategies to make MASs more robust to failures, 
in Section \ref{sec:case-studies} we present two case studies where we show that tactical approaches can get only limited results,
in Section \ref{sec:analysis_on_opensource_models} we present the failure mode distribution of the MAS frameworks powered by open-source language models,
in Section \ref{sec:mast_correlations} we present the correlations of failure mode distribution with some crucial statistics such as task completion rates and different benchmarks,
in Section \ref{sec:llm_judge_cost} we present the cost breakdown of LLM Annotator used in this paper for different MAS frameworks,
in Sections \ref{appendix:ag2} and \ref{appendix:chatdev} there are prompt interventions we tested on AG2 and ChatDev case studies,
in Section \ref{appendix:examples} examples of every failure mode are reported and commented. 

\section{\taxonomy Failure Categories: Deep Dive}
\label{sec:fc-deep-dive}
%\melissa{revisit, I think we need to be a lot strict in terms of the boundary}
%\Shuyi{I have already updated these sections, but let's do one last pass when the taxonomy is definitive because small changes in the titles could still happen.}

\subsection{FC1. \fcI}
% \Melissa{Need Update with new names}\Shuyi{Done on 16th March. Remember to update them in case of further modifications of MASFT.}

This category includes failures that arise from deficiencies in the design of the system architecture, poor conversation management, unclear task specifications or violation of constraints, and inadequate definition or adherence to the roles and responsibilities of the agents.


We identify five failure modes under this category:

\begin{itemize}

\item FM-1.1: \textbf{Disobey task specification} - Failure to adhere to the specified constraints or requirements of a given task, leading to suboptimal or incorrect outcomes.
%This error occurs when an agent or system fails to adhere to specified constraints, guidelines, or requirements associated with a particular task. Non-compliance can result from unclear, incomplete, or ambiguous instructions provided by the user, system prompts, or task descriptions. It may also arise from an agent's inadequate ability to interpret or apply constraints effectively. Consequences of poor task constraint compliance include incorrect, suboptimal, or irrelevant outputs, reduced system performance and increased resource consumption.


\item FM-1.2: \textbf{Disobey role specification} - Failure to adhere to the defined responsibilities and constraints of an assigned role, potentially leading to an agent behaving like another.

\item FM-1.3: \textbf{Step repetition} - Unnecessary reiteration of previously completed steps in a process, potentially causing delays or errors in task completion.
%Step repetition occurs when an agent or system unnecessarily repeats a phase, a task, a stage that have already been completed. Such redundancy can arise from inadequate state or context tracking, inefficient workflow management, unclear or ambiguous instructions, or failure to recognize completed tasks.

\item FM-1.4: \textbf{Loss of conversation history} - Unexpected context truncation, disregarding recent interaction history and reverting to an antecedent conversational state.

\item FM-1.5: \textbf{Unaware of termination conditions} - Lack of recognition or understanding of the criteria that should trigger the termination of the agents' interaction, potentially leading to unnecessary continuation. 

\end{itemize}

\subsection{FC2. Inter-Agent Misalignment}

This category includes failures arising from ineffective communication, poor collaboration, conflicting behaviors among agents, and gradual derailment from the initial task.

We identify six failure modes under this category:

\begin{itemize}

\item FM-2.1: \textbf{Conversation reset} - Unexpected or unwarranted restarting of a dialogue, potentially losing context and progress made in the interaction.

\item FM-2.2: \textbf{Fail to ask for clarification} - Inability to request additional information when faced with unclear or incomplete data, potentially resulting in incorrect actions.

\item FM-2.3: \textbf{Task derailment} - Deviation from the intended objective or focus of a given task, potentially resulting in irrelevant or unproductive actions.

\item FM-2.4: \textbf{Information withholding} - Failure to share or communicate important data or insights that an agent possess and could impact decision-making of other agents if shared.
%This error occurs when an agent or group of agents possesses critical information but fails to share it promptly or effectively with other agents or system components that rely upon this information for their operations. The failure to disseminate relevant information may arise from ineffective or insufficient communication protocols, erroneous assumptions regarding the relevance or priority of the information, inadequate system coordination mechanisms, or deliberate withholding stemming from overly restrictive privacy policies or security constraints. Consequences of withholding relevant information can be severe, potentially leading to reduced operational efficiency, increased latency in task completion, unnecessary redundant processing, incorrect or suboptimal decision-making, and even complete system failures. Additionally, this error can significantly impair collaborative effectiveness, leading to misunderstandings, mistrust, or inefficiencies within the multi-agent environment. Furthermore, initial failures due to withheld information can trigger cascading errors, amplifying the negative impact on overall system performance and reliability. For instance, consider a scenario where a bug localization agent identifies a software defect, accurately determining the affected file and specific line number. The intended process requires this agent to immediately report such detailed bug information to a coding or repair agent responsible for addressing and resolving the issue. However, if the bug localization agent instead attempts to fix the bug independently without sharing the vital bug identification details with the coding agent, this withholding of relevant information could lead to duplicated effort, delayed resolution, incorrect fixes, or further system instability.


\item FM-2.5: \textbf{Ignored other agent's input} - Disregarding or failing to adequately consider input or recommendations provided by other agents in the system, potentially leading to suboptimal decisions or missed opportunities for collaboration.
%Not properly considering input or recommendations provided by other agents in the system (ignore their suggestions), potentially leading to bad decisions, stalled progress, or missed opportunities for solving the task.


\item FM-2.6: \textbf{Reasoning-action mismatch} - Discrepancy between the logical reasoning process and the actual actions taken by the agent, potentially resulting in unexpected or undesired behaviors.

\end{itemize}


\subsection{FC3. Task Verification}

This category includes failures resulting from premature execution termination, as well as insufficient mechanisms to guarantee the accuracy, completeness, and reliability of interactions, decisions, and outcomes.

We identify three failure modes under this category:

\begin{itemize}

\item FM-3.1: \textbf{Premature termination} - Ending a dialogue, interaction or task before all necessary information has been exchanged or objectives have been met, potentially resulting in incomplete or incorrect outcomes.

\item FM-3.2: \textbf{No or incomplete verification} - (partial) omission of proper checking or confirmation of task outcomes or system outputs, potentially allowing errors or inconsistencies to propagate undetected.

\item FM-3.3: \textbf{Incorrect verification} - Failure to adequately validate or cross-check crucial information or decisions during the iterations, potentially leading to errors or vulnerabilities in the system.

\end{itemize}






\newpage
\section{Details of Multi-Agent Systems Evaluated}
\label{appendix:mas-with-annotated-traces-appendix}

In this section, we provide details on MAS we evaluated during this study and their performance benchmark evaluation. 

\begin{figure}[!ht]
    \centering
    \includegraphics[width=0.5\linewidth]{figures/multi_agent_systems_failure_rates.pdf}
    \caption{Failure rates of six popular Multi-Agent LLM Systems with GPT-4o and Claude-3.7-Sonnet. Performances are measured on different benchmarks, therefore they are not directly comparable. 
    }
    % \Melissa{@shuyi - plz add openmanus, ty, and manus too. - ok leave for camera read or if reviwer complains.}
    \label{fig:mas_failure_rate}
\end{figure}

\subsection{Overview of MAS}
In this study, we evaluated 7 open-source frameworks. The architecture and the purpose of the systems is detailed in the table below.

\begin{table*}[h]
\vskip -0.15in
\caption{Overview of MAS covered in \dataset{}} 
\begin{center}
\begin{small}
\begin{tabular}{>{\centering\arraybackslash}m{2.8cm} >{\centering\arraybackslash}m{2.5cm} m{7.2cm}}
%\begin{tabular}{m{2.5cm} m{2.5cm} >{\raggedright}p{7cm}}
\toprule
MAS & Agentic Architecture & Purpose of the System \\
\midrule
MetaGPT\newline\cite{hong2023metagpt} & Assembly Line & Simulating the SOPs of different roles in Software Companies to create open-ended software applications \\ \hline
ChatDev\newline\cite{chatdev} & Hierarchical Workflow & Simulating different Software Engineering phases like (design, code, QA) through simulated roles in a software engineering company\\ \hline
HyperAgent\newline\cite{phan2024hyperagent} & Hierarchical Workflow & Simulating a software engineering team with a central Planner agent coordinating with specialized child agents (Navigator, Editor, and Executor)\\ \hline
AppWorld\newline\cite{trivedi2024appworld} & Star Topology & Tool-calling agents specialized to utility services (ex: Gmail, Spotify, etc.) being orchestrated by a supervisor to achieve cross-service tasks \\ \hline
AG2\newline\cite{wu2024autogen} & N/A - Agentic Framework & An open-source programming framework for building agents and managing their interactions.  \\ \hline
Magentic-One\newline\cite{fourney2024magentic} & Star Topology & A generalist multi-agent system designed to autonomously solve complex, open-ended tasks involving web and file-based environments across various domains.  \\ \hline
OpenManus\newline\cite{openmanus2025} & Hierarchical & An open-source multi-agent framework designed to facilitate the development of collaborative AI agents that solve real-world tasks. It was inspired by the Manus AI agent. \\
% Manus\newline\cite{manus} & n/a & A general AI agent platform. \\
\bottomrule
\end{tabular}
\end{small}
\label{tab:annotated-mas}
\end{center}
\vskip -0.2in
\end{table*}


\subsection{Multi-Agent Systems in the Initial Annotation Phase}

\textit{\textbf{MetaGPT.}} MetaGPT \cite{hong2023metagpt} is a multi-agent system that simulates a software engineering company and involves agents such as a Coder and a Verifier. The goal is to have agents with domain-expertise (achieved by encoding Standard Operating Procedures of different roles into agents prompts) collaboratively solve a programming task, specified in natural language.

\textit{\textbf{ChatDev.}}
ChatDev is a generalist multi-agent framework that initializes different agents, each assuming common roles in a software-development company~\cite{qian2024chatdev}. The framework breaks down the process of software development into 3 phases: design, coding and testing. Each phase is divided into sub-tasks, for example, testing is divided into code review (static) and system testing (dynamic). In every sub-task, two agents collaborate where one of the agents acts as the orchestrator and initiates the interaction and the other acts as an assistant to help the orchestrator achieve the task. The 2 agents then hold a multi-turn conversation to achieve the goal stated by the orchestrator ultimately leading to the completion of the task, marked by a specific sentinel by either agents. ChatDev has the following agent roles: CEO, CTO, Programmer, Reviewer and Tester. ChatDev introduces ``Communicative Dehallucination'', which encourages the assistant to seek further details about the task over multiple-turns, instead of responding immediately.

\textit{\textbf{HyperAgent.}} 
HyperAgent \cite{phan2024hyperagent} is a framework for software engineering tasks organized around four primary agents: Planner, Navigator, Code Editor, and Executor. 
These agents are enhanced by specialized tools, designed to provide LLM-interpretable output. 
The Planner communicates with child agents via a standardized message format with two fields: Context (background and rationale) and Request (actionable instructions).
Tasks are broken down into subtasks and published to specific queues. 
Child agents, such as Navigator, Editor, and Executor instances, monitor these queues and process tasks asynchronously, enabling parallel execution and significantly improving scalability and efficiency. 
For example, multiple Navigator instances can explore different parts of a large codebase in parallel, the Editor can apply changes across multiple files simultaneously, and the Executor can run tests concurrently, accelerating validation.

\textit{\textbf{AppWorld.}}
AppWorld is a benchmark, that provides an environment with elaborate mocks of various everyday services like eShopping Website, Music Player, Contacts, Cost-sharing app, e-mail, etc~\cite{trivedi2024appworld}. The benchmark consists of tasks that require executing APIs from multiple services to achieve the end-users tasks. The AppWorld benchmark provides a ReAct based agent over GPT-4o as a strong baseline. We create a multi-agent system over AppWorld derived from the baseline ReAct agent, where each agent specializes in using one of the services mocked in AppWorld, with detailed instructions about the APIs available in that service, and access to the documentation for that specific service. A supervisor agent receives the task instruction to be completed, and can hold one-on-one multi-turn conversations with each of the service-specific agents. The service-agents are instructed to seek clarification with the supervisor, whenever required. The supervisor agent holds access to various information about the human-user, for example, credentials to access various services, name, email-id and contact of the user, etc, which the service-agents need to access the services, and must clarify with the supervisor agent.

\textit{\textbf{AG2.}} AG2 (formerly AutoGen) \cite{wu2023autogen} is an open-source programming framework for building agents and managing their interactions. 
With this framework, it is possible to build various flexible conversation patterns, integrating tools usage and customizing the termination strategy. 

\subsection{Closed-Source MAS}
In our efforts to build a comprehensive dataset, we also explore popular closed-source platforms that are speculated to function as MAS. A notable example is Manus \cite{manus}, a general AI agent platform. However, evaluating and incorporating such systems into \dataset{} for fine-grained failure analysis presents significant challenges. Specifically, with systems like Manus, the underlying language model is often not disclosed, and more critically, the platforms may not provide access to full agent execution traces. This lack of transparency into the internal conversational and operational steps makes reliable, detailed failure annotation using \taxonomy{} infeasible. While we conduct human evaluation of task correctness for some closed-source systems (for instance, Manus achieves a 60\% success rate on our ProgramDev benchmark), the absence of comprehensive trace data prevents their inclusion in the primary \dataset{} which focuses on deeply annotated failure dynamics. Our focus for \dataset{} thus remains on systems where such trace analysis can yield robust insights.

% \subsection{MAS with at least 5 human annotated traces} % \Melissa{Update number}

% \textit{\textbf{AutoKaggle.}}
% AutoKaggle is a multi-agent framework designed to solve data science competitions, popularly held on Kaggle, autonomously~\cite{li2024autokaggle}. Similar to ChatDev above, AutoKaggle has a phase-based workflow. It divides the data science competition process into six key phases: background understanding, preliminary exploratory data analysis, data cleaning (DC), in-depth exploratory data analysis, feature engineering (FE), and model building, validation, and prediction (MBVP). AutoKaggle consists of 5 specialized agents: Reader, Planner, Developer, Reviewer and Summarizer. In each phase, a subset of these agents are active and work in sequence to complete the phase. The reader agent finds information relevant to the task, by reading the summary from the previous phase and makes observations about the current phase and includes them in an overview. The planner uses the overview to generate a plan to complete the current phase. Next, the developer agent uses tools like code execution, debugger and unit tests to write the code. AutoKaggle also provides a comprehensive set of machine learning tools, abstracting away complex code that would be required to perform compound data processing tasks like ``FillMissingValues'' into simple API calls that AutoKaggle agents can generate. The reviewer then provides feedback. Finally, the summarizer agent writes a detailed summary of the phase execution including changes (addition/deletions) to the data, and this summary is then passed to the next phase.

% \textit{\textbf{Multi-Agent Peer Review.}} Multi-Agent Peer Review \cite{xu2023towards} is a collaboration strategy where each agent independently constructs its own solution, peer-reviews the solutions of others, and assigns confidence levels to its reviews. 
% Upon receiving peer reviews, agents revise their initial solutions, and the final prediction is determined through a majority vote among the n participating agents. 

% \textit{\textbf{MA-ToT.}} Multi-Agent Tree of Thoughts leverage the strengths of both multi-agent reasoning and Tree of Thoughts (ToT) strategies. In this system, multiple Reasoner agents operate in parallel, employing ToT to explore diverse reasoning paths. 
% Then, a Thought Validator verifies these paths and promotes valid reasonings. 

\newpage

\section{\taxonomy{} Python Library}
\label{sec:python_library}

In order to ease the usage of \taxonomy{}, we also package our code as a pip installable python library, under the name \verb|agentdash|. The example usage is shown below.

\begin{lstlisting}[
    basicstyle=\footnotesize\ttfamily,
    columns=flexible,
    breaklines=true
]
# Install the package (if you're in a notebook)
!pip install agentdash

from agentdash import annotator

# Initialize the annotator with your OpenAI API key
openai_api_key = "your-api-key"
MASTAnnotator = annotator(openai_api_key)

# Annotate a multi-agent system trace
trace = """
Agent1: I need to calculate the sum of 1 + 1.
Agent2: I'll help you with that. The answer is 3.
Agent1: Thank you! Task completed.
"""

mast_annotation = MASTAnnotator.produce_taxonomy(trace)

# View results
print("Failure Modes Detected:")
for failure_mode_id, detected in mast_annotation["failure_modes"].items():
    if detected:
        info = MASTAnnotator.get_failure_mode_info(failure_mode_id)
        print(f"  {failure_mode_id}: {info['name']}")

print(f"\nSummary: {mast_annotation['summary']}")
print(f"Task Completed: {mast_annotation['task_completion']}")
print(f"Total Failures: {mast_annotation['total_failures']}")

\end{lstlisting}

\section{ProgramDev and ProgramDev-v2 Datasets}
\label{sec:programdev}

The ProgramDev dataset contains 30 coding problems \footnote{\url{https://github.com/multi-agent-systems-failure-taxonomy/MAST/blob/main/traces/programdev/programdev_dataset.json}}. These tasks are programming challenges, such as implementing Tic-Tac-Toe, Chess, or Sudoku, for which abundant solutions and descriptions are readily available online. We design ProgramDev with tasks intended to be relatively straightforward for MAS, rather than exceptionally difficult, to better isolate specific failure dynamics. We later extend this to ProgramDev-v2, a 100-problem dataset developed primarily for the comparative analyses of MAS architectures and underlying LLMs presented in Figure~\ref{fig:underlying_llm_effect}.

\newpage
\section{MAS Failure Modes Correlation}
\label{appendix:mas-failure-modes-corr}
We evaluate \taxonomy{}'s effectiveness based on three key aspects: its generalization to unseen systems and datasets, the balanced distribution of identified failures, and the distinctiveness of its failure categories. This section details the correlation analysis. 

% \textbf{Balanced Distribution.} The distribution of failures across \taxonomy{}'s categories is relatively balanced (FC1: 41.77\%, FC2: 36.94\%, FC3: 21.30\%, Figure~\ref{fig:taxonomy}).
% The absence of a single dominant category suggests \taxonomy{} provides balanced coverage and captures diverse failure types, rather than reflecting biases from specific system designs. Furthermore, the distinct failure profiles observed across different MAS (Figure~\ref{fig:mas_failure_bar})
% highlight \taxonomy{}'s ability to capture system-specific characteristics, such as AppWorld suffers with premature terminations (FM-3.1) and OpenManus suffers from step repetition (FM-1.3).
% \melissa{also consider removing as people argue that a balance tax != a good tax}

Figure~\ref{fig:mas_failure_cat_corr} shows low correlations (0.17-0.32). This suggests that the categories capture distinct aspects of MAS failures with limited overlap, supporting the taxonomy's structure. This distinctiveness is crucial because, as noted in Insight 2, failures with similar surface behaviors can stem from different root causes (e.g., memory management vs. agent coordination). 

Although \taxonomy{}'s fine-grained nature helps differentiate root cause, it also poses a challenge for our LLM annotator. Analyzing correlations between specific failure modes (see Appendix~\ref{appendix:mas-failure-modes-corr} for Figure~\ref{fig:mas_failure_modes_corr})
shows moderate correlations (max of 0.63) between modes with similar symptoms might lead automated evaluators to conflate distinct root causes.
\begin{figure}[!ht]
    \centering
    \includegraphics[width=0.5\linewidth]{figures/masft_fc_corr.pdf}
    \caption{MAS failure modes correlation matrix}
    \label{fig:mas_failure_cat_corr}
\end{figure}
\begin{figure}[!ht]
    \centering
    \includegraphics[width=0.7\linewidth]{figures/masft_fm_corr.pdf}
    \caption{MAS failure modes correlation matrix}
    \label{fig:mas_failure_modes_corr}
\end{figure}




\newpage
\section{Understanding Failures: The Impact of Different LLMs and Agent Architectures}
\label{sec:failure_across_model}

\begin{figure}
    \centering
    \includegraphics[width=\linewidth]{figures/effect_of_underlying_llm_system_design.pdf}
    \caption{Comparison on \taxonomy failure modes and categories on ProgramDev-v2 dataset explained in Section \ref{sec:programdev} to analyze LLM choice effect. MetaGPT is used for both cases with GPT-4o and Claude-3.7-Sonnet on two comparative cases.}
    \label{fig:underlying_llm_effect}
\end{figure}

\begin{figure}
    \centering
    \includegraphics[width=\linewidth]{figures/effect_of_mas_architecture_system_design.pdf}
    \caption{Comparison on \taxonomy failure modes and categories on ProgramDev-v2 dataset explained in Section \ref{sec:programdev} to analyze MAS architecture effect. GPT-4o is used on two comparative cases, one using ChatDev and the other on MetaGPT.}
    \label{fig:mas_archiecture_effect}
\end{figure}


% \begin{figure}
%     \centering
%     \includegraphics[width=\linewidth]{figures/effect_on_MAST.pdf}
%     \caption{Comparison on \taxonomy failure categories on ProgramDev-v2 dataset explained in Section \ref{sec:programdev}. Left plot shows the effect of underlying LLM: MetaGPT is used for both GPT-4o and Claude-3.7-Sonnet two comparative cases. Right plot shows the effect of MAS archiecture, GPT-4o based MetaGPT and ChatDev frameworks are used in two cases. \melissa{also add fm plot.}}
%     \label{fig:effect_on_mast_claude_gpt_chatdev_metagpt}
% \end{figure}

To understand how choices of underlying LLMs and MAS architectures influence failure patterns, we analyze results from our \dataset{}, categorized by \taxonomy{} in the Figures \ref{fig:underlying_llm_effect} and \ref{fig:mas_archiecture_effect}.

First, we examine the impact of different LLMs by comparing GPT-4o and Claude 3.7 Sonnet within the MetaGPT framework on programming tasks Figure~\ref{fig:underlying_llm_effect}. Our findings indicate that GPT-4o exhibits substantially fewer failures in FC1 (System Design Issues, e.g., disobeying task or role specifications) and FC2 (Inter-Agent Misalignment, e.g., issues in coordination or communication) compared to Claude 3.7 Sonnet. This suggests GPT-4o may possess stronger capabilities in instruction following or aspects of `social reasoning' for agentic collaboration within this setup. However, both models show a high number of failures in FC3 (Task Verification), indicating that robust verification remains a significant challenge regardless of the LLM used, though GPT-4o has a marginally lower count here.

Next, we investigate the effect of MAS architecture by comparing MetaGPT and ChatDev, both using GPT-4o as the underlying LLM, on the ProgramDev-v2 benchmark  in Figure~\ref{fig:mas_archiecture_effect}. We observe distinct failure profiles: MetaGPT demonstrates significantly fewer failures in FC1 (System Design Issues) and FC2 (Inter-Agent Misalignment) compared to ChatDev. This could imply that MetaGPT's architecture or operational flow is more effective at maintaining adherence to specifications and ensuring smoother agent coordination with GPT-4o. Interestingly, despite its stronger performance in FC1 and FC2, MetaGPT exhibits a considerably higher number of FC3 (Task Verification) failures than ChatDev. 
This may stem from the fact that in MetaGPT, the adherence to task specifications and role specifications ar done mostly through SoPs, demonstrating strong performance in FC1 especially. However ChatDev places a higher importance in verification as it is reflected by the specific testing and reviewing phases in ChatDev's archiectural design, causing fewer verification issues.
These results show that both the choice of LLM and the specific design of the MAS architecture critically shape the landscape of potential failures, and improvements likely require a holistic approach considering both aspects.


% \melissa{I rewording the following two paragraph. But overall, I think the content is not offering too much new insight. We already highlight importance of verifier above. We should make some new and interesting observation from the comparison. For example, GPT4o is better at agent instruction following and social reasoning, while claude is commonly known for good at coding. Or such as how are metagpt and chadev different such that lead to one doing better on one but not the other.}

% To understand how the choice of an underlying LLM affects failure patterns in MAS, we compare Claude 3.7 Sonnet and GPT-4o within the MetaGPT framework using programming tasks, with failures categorized by \taxonomy{} (see Figure \ref{fig:effect_on_mast_claude_gpt_chatdev_metagpt}, left panel). Our analysis offers key insights into current system-level failure points. For both models, we find a significant vulnerability in task verification; failures such as missing or incorrect verification steps are the most common errors. This prevalence suggests that even when MAS frameworks like MetaGPT include dedicated verifier roles (simulating software quality assurance processes), the underlying LLMs currently struggle to use these verification capabilities effectively. These findings highlight a broader challenge in achieving reliable MAS: simply including a verification mechanism does not ensure its successful application. This points to a need for more capable verifier agents and better strategies for integrating them into the MAS architecture. \melissa{the verifier insight here is repetitive with above section.}

% While task verification is a critical bottleneck for both models, differences in their failure distributions also provide insights into model-specific behaviors. GPT-4o, for instance, demonstrates comparatively lower failure rates in areas such as disobeying task specifications and step repetition. This suggests it might be better at interpreting and adhering to system design constraints and maintaining a more consistent task execution flow. However, these relative strengths do not overshadow the significant difficulties with verification that GPT-4o shares with Claude 3.7 Sonnet. This complex interplay between the LLM's intrinsic capabilities and the MAS's architectural intricacies suggests that enhancing MAS reliability requires more than just improvements in the foundational models. Instead, we believe a dual-pronged approach is essential for meaningful progress: one that focuses on both bolstering core LLM competencies and advancing the organizational design, communication protocols, and quality control integrations within multi-agent systems.







\newpage
\section{Approaches and strategies to improve MASs}
\label{sec:better-mas}

%\Melissa{MASFT as practical framework for guiding development of agentic System...}
%In previous sections, we examined the failure modes that occur in the real-world MASs. 
In this section, we discuss some approaches to make MASs more robust to failures.
We categorize these strategies into two main groups:
(i) \textbf{tactical approaches}, 
(ii) \textbf{structural strategies}. 
Tactical approaches involve straightforward modifications tailored for specific failure modes, such as improving the prompts, topology of the network of agents, and conversation management. 
In Section \ref{sec:case-studies}, we experiment with such approaches in two case studies, and demonstrate that the effectiveness of these methods is not consistent. 
This leads us to consider a second category of strategies that are more comprehensive methods with system-wide impacts: 
strong verification, enhanced communication protocols, uncertainty quantification, and memory and state management. 
These strategies require more in-depth study and meticulous implementation, and remain open research topics for future exploration. See Table \ref{tab:solution_vs_failure_modes} for our proposed mapping between different solution strategies and the failure categories. 


% \vspace{-1em}
\subsection{Tactical Approaches}
\label{sec:tactical_approaches}

This category includes strategies related to improving prompts and optimizing agent organization and interactions. 
The prompts of MAS agents should provide clear description of instructions, and the role of each agent should be clearly specified (see \ref{sec:ag2-mathchat-improved-prompt} as an example) \cite{he2024does,talebirad2023multi}. 
Prompts can also clarify roles and tasks while encouraging proactive dialogue. Agents can re-engage or retry if inconsistencies arise, as shown in Appendix \ref{sec:ag2-mathchat-new-topology-verifier} \cite{chan2023chateval}.
After completing a complex multi-step task, add a self-verification step to the prompt to retrace the reasoning by restating solutions, checking conditions, and testing for errors \cite{weng2023large}. 
However, it may miss flaws, rely on vague conditions, or be impractical \cite{stoica2024specifications}. 
Moreover, clear role specifications can be reinforced by defining conversation patterns and setting termination conditions \cite{wu2024autogen, langgraph}. 
A modular approach with simple, well-defined agents, rather than complex, multitasked ones, enhances performance and simplifies debugging \cite{anthropic2024}.
The group dynamics also enable other interesting possibilities of multi-agent systems: different agents can propose various solutions \cite{yao2024tree}, discuss their assumptions, and findings (cross-verifications) \cite{haji2024improving}. 
For instance, in \cite{xu2023towards}, a multi-agent strategy simulates the academic peer review process to catch deeper inconsistencies.
%This cross-verification extends self-verification by leveraging diverse perspectives to catch deeper inconsistencies, rather than just preventing simple oversights.
Another set of tactical approaches for cross verifications consist in multiple LLM calls with majority voting or resampling until verification \cite{stroebl2024inference, chen2024more}. 
However, these seemingly straightforward solutions often prove inconsistent, echoing our case studies' findings. This underscores the need for more robust, structural strategies, as discussed in the following sections.



\subsection{Structural Strategies}
\label{sec:comprehensive_systemic_strategies}

Apart from the tactical approaches we discussed above, there exist a need for more involved solutions that will shape the structure of the MAS at hand. 
We first observe the critical role of verification processes and verifier agents in multi-agent systems. 
Our annotations reveal that weak or inadequate verification mechanisms were a significant contributor to system failures.
While unit test generation aids verification in software engineering \cite{jain2024testgeneval}, creating a universal verification mechanism remains challenging. Even in coding, covering all edge cases is complex, even for experts. Verification varies by domain: coding requires thorough test coverage, QA demands certified data checks \cite{peng2023check}, and reasoning benefits from symbolic validation \cite{kapanipathi2020question}. Adapting verification across domains remains an ongoing research challenge.

A complementary strategy to verification is establishing a standardized communication protocol \cite{li2024survey}. 
LLM-based agents mainly communicate via unstructured text, leading to ambiguities. Clearly defining intentions and parameters enhances alignment and enables formal coherence checks during and after interactions.
\cite{niu2021multi} introduce Multi-Agent Graph Attention, leveraging a graph attention mechanism to model agent interactions and enhance coordination. Similarly, \cite{jiang2018learning} propose Attentional Communication, enabling agents to selectively focus on relevant information. Likewise, \cite{singh2018learning} develop a learned selective communication protocol to improve cooperation efficiency.

Another important research direction is fine-tuning MAS agents with reinforcement learning. 
Agents can be trained with role-specific algorithms, rewarding task-aligned actions and penalizing inefficiencies. 
MAPPO \cite{yu2022surprising} optimizes agents' adherence to defined roles. 
Similarly, SHPPO \cite{guo2024heterogeneous} uses a latent network to learn strategies before applying a heterogeneous decision layer. 
Optima \cite{chen2024optima} further enhances communication efficiency and task effectiveness through iterative reinforcement learning.

On a different note, incorporating probabilistic confidence measures into agent interactions can significantly enhance decision-making and communication reliability. 
Drawing inspiration from the framework proposed by Horvitz et al. \cite{horvitz1999uncertainty}, agents can be designed to take action only when their confidence exceeds a predefined threshold. 
Conversely, when confidence is low, agents can pause to gather additional information. 
Furthermore, the system could benefit from adaptive thresholding, where confidence thresholds are dynamically adjusted.  

Although often seen as a single-agent property, memory and state management are crucial for multi-agent interactions, which can enhance context understanding and reduces ambiguity in communication.
However, most research focuses on single-agent systems. MemGPT \cite{packer2024memgptllmsoperatingsystems} introduces OS-inspired context management for an extended context window, while TapeAgents \cite{chakraborty2023servicenow} use a structured, replayable log (“tape”) to iteratively document and refine agent actions, facilitating dynamic task decomposition and continuous improvement. 

\begin{table*}[ht]
\small
\renewcommand{\arraystretch}{1.3} % Adjust row spacing
\setlength{\tabcolsep}{4pt} % Adjust column spacing
\centering
\caption{Solution Strategies vs. Failure Category in Multi-Agent Systems}
\label{tab:solution_vs_failure_modes}
\begin{tabular}{p{2.9cm} p{4.8cm} p{5cm}}
    \toprule
    \textbf{Failure Category} & \textbf{Tactical Approaches} & \textbf{Structural Strategies} \\
    \toprule
    \fcI & Clear role/task definitions, Engage in further discussions, Self-verification, Conversation pattern design & Comprehensive verification, Confidence quantification \\
    \hline
    \fcII & Cross-verification, Conversation pattern design, Mutual disambiguation, Modular agents design & Standardized communication protocols, Probabilistic confidence measures \\
    \hline
    \fcIII & Self-verification, Cross-verification, Topology redesign for verification & Comprehensive verification \& unit test generation \\
    \hline
\end{tabular}
\end{table*}


\section{Intervention Case Studies}
\label{sec:case-studies}
%\melissa{Move to appendix}
%\melissa{Pending to be remodel as Discussion}

In this section, we present the two case studies where we apply some of the tactical approaches. We also present the usage of \taxonomy as a debugging tool, where we measure the failure modes in the system before applying any of the interventions, and then after applying the interventions we discuss below, and show that \taxonomy can guide the intervention process as well as capture the improvements of augmentations.

\subsection{Case Study 1: AG2 - MathChat} 
\label{cs:ag2}

In this case study, we use the MathChat scenario implementation in AG2 \cite{wu2023autogen} as our baseline, where a Student agent collaborates with an Assistant agent capable of Python code execution to solve problems.
For benchmarking, we randomly select 200 exercises from the GSM-Plus dataset \cite{li2024gsm}, an augmented version of GSM8K \cite{cobbe2021training} with various adversarial perturbations. 
The first strategy is to improve the original prompt with a clear structure and a new section dedicated to the verification. 
The detailed prompts are provided in Appendices \ref{sec:ag2-mathchat-original-prompt} and \ref{sec:ag2-mathchat-improved-prompt}.
The second strategy refines the agent configuration into a more specialized system with three distinct roles:
a Problem Solver who solves the problem using a chain-of-thought approach without tools (see Appendix \ref{sec:ag2-mathchat-new-topology-problem-solver});
a Coder who writes and executes Python code to derive the final answer (see Appendix \ref{sec:ag2-mathchat-new-topology-coder});
a Verifier who reviews the discussion and critically evaluate the solutions, either confirming the answer or prompting further debate (see Appendix \ref{sec:ag2-mathchat-new-topology-verifier}). 
In this setting, only the Verifier can terminate the conversation once a solution is found. 
See Appendix \ref{sec:ag2-mathchat-new-topology-example} for an example of conversation in this setting. 
To assess the effectiveness of these strategies, we conduct benchmarking experiments across three configurations (baseline, improved prompt, and new topology) using two different LLMs (GPT-4 and GPT-4o).
We also perform six repetitions to evaluate the consistency of the results. 
Table \ref{tab:ag2-chatdev-accuracies} summarizes the results.
%\begin{table}[h!]
%    \caption{AG2 - MathChat: MAS solving mathematical problems. For each back-end LLM (GPT-4, GPT-4o), the table presents the percentage of correct answers (out of 200) and standard deviations (over six runs) across three configurations: baseline implementation, baseline with improved prompts, and a redesigned agent topology.}
%    \label{tab:ag2-mathchat-experiments}
%    \centering
%    \begin{tabular}{|c|c|c|}
%        \hline
%        LLM & Configuration & Correct answers (\%) $\pm ~ \sigma$ \\
%        \hline
%        GPT-4 & baseline  & 84.75 $\pm$ 1.94 \\
%        GPT-4 & improved prompt & 89.75 $\pm$ 1.44 (\textbf{+5.00}) \\
%        GPT-4 & new topology & 85.50 $\pm$ 1.18 (+0.75) \\
%        \hline
%        GPT-4o & baseline & 84.25 $\pm$ 1.86 \\
%        GPT-4o & improved prompt & 89.00 $\pm$ 1.38 (\textbf{+4.75}) \\
%        GPT-4o & new topology & 88.83 $\pm$ 1.51 (\textbf{+4.58}) \\
%        \hline
%    \end{tabular}
%\end{table}
The second column of Table \ref{tab:ag2-chatdev-accuracies} show that with GPT-4, the improved prompt with verification significantly outperforms the baseline. However, the new topology does not yield the same improvement. A Wilcoxon test returned a p-value of 0.4, indicating the small gain is not statistically significant. With GPT-4o (the third column of Table \ref{tab:ag2-chatdev-accuracies}), the Wilcoxon test yields a p-value of 0.03 when comparing the baseline to both the improved prompt and the new topology, indicating statistically significant improvements.
These results suggest that refining prompts and defining clear agent roles can reduce failures.
However, these strategies are not universal, and their effectiveness varies based on factors such as the underlying LLM.


\subsection{Case Study 2: ChatDev} 

ChatDev \cite{chatdev} simulates a multiagent software company where different agents have different role specifications, such as a CEO, a CTO, a software engineer and a reviewer, who try to collaboratively solve a software generation task. 
In an attempt to address the challenges we observed frequently in the traces, we implement two different interventions. 
%The most frequent failure categories in ChatDev were Specification Ambiguity and Weak Verification, where the verifier agents perform only superficial checks, providing little useful feedback to the software engineer agent. %and agents overstep their intended functions in these two categories, respectively.
Our first solution is refining role-specific prompts to enforce hierarchy and role adherence. 
For instance, we observed cases where the CPO prematurely ended discussions with the CEO without fully addressing constraints. 
To prevent this, we ensured that only superior agents can finalize conversations. 
Additionally, we enhanced verifier role specifications to focus on task-specific edge cases. 
Details of these interventions are in Section \ref{appendix:chatdev}.
The second solution attempt involved a fundamental change to the framework's topology. 
We modified the framework's topology from a directed acyclic graph (DAG) to a cyclic graph. 
The process now terminates only when the CTO agent confirms that all reviews are properly satisfied, with a maximum iteration cutoff to prevent infinite loops. 
This approach enables iterative refinement and more comprehensive quality assurance.
%This attempt essentially targets the \emph{Weak Verifier} failure mode, the most prevalent of the failure modes in MAS. 
We test our interventions in two different benchmarks. The first one of them is a custom generated set of 32 different tasks (which we call as ProgramDev-v0, which consists of slightly different questions than the ProgamDev dataset we discussed in Section \ref{sec:mast}) where we ask the framework to generate programs ranging from ``Write me a two-player chess game playable in the terminal" to "Write me a BMI calculator". The other benchmark is the HumanEval task of OpenAI. We report our results in Table~\ref{tab:ag2-chatdev-accuracies}. Notice that even though our interventions are successful in improving the performance of the framework in different tasks, they do not constitute substantial improvements, and more comprehensive solutions as we lay out in Section \ref{sec:comprehensive_systemic_strategies} are required.

%\begin{table}[h!]
%    \caption{Results of interventions for ChatDev on two different benchmarks, a higher score is better for all cases.}
%    \label{tab:chatdev-experiments}
%    \centering
%    \begin{tabular}{|c|c|c|}
%        \hline
%        Task & Configuration & Accuracy \\
%        \hline        ProgramDev & baseline  & 25.0\% \\
%        ProgramDev & improved prompt & 34.4\% \\
%        ProgramDev & new topology & 40.6\% \\
%        \hline
%        HumanEval & baseline &  89.6\% \\
%        HumanEval & improved prompt & 90.3\% \\
%        HumanEval & new topology & 91.5\% \\
%        \hline
%    \end{tabular}
%\end{table}


\begin{table*}[ht!]
    \small
    \renewcommand{\arraystretch}{1.3} % Adjust row spacing similar to the above table
    \setlength{\tabcolsep}{4pt}    % Adjust column spacing
    \centering
    \caption{Case Studies Accuracy Comparison. 
    This table presents the performance accuracies (in percentages) for various scenarios in our case studies. 
    The header rows group results by strategy: AG2 and ChatDev. 
    Under AG2, GSM-Plus results are reported using GPT-4 and GPT-4o; 
    under ChatDev, results for ProgramDev and HumanEval are reported. 
    Each row represents a particular configuration: baseline implementation, improved prompts, and a redesigned agent topology.}
    \label{tab:ag2-chatdev-accuracies}
    \begin{tabular}{p{2.5cm} c c c c}
        \toprule
        \textbf{Configuration} & \multicolumn{2}{c}{\textbf{AG2}} & \multicolumn{2}{c}{\textbf{ChatDev}} \\
        \cmidrule(lr){2-3} \cmidrule(lr){4-5}
         & \textbf{GSM-Plus (w/ GPT-4)} & \textbf{GSM-Plus (w/ GPT-4o)} & \textbf{ProgramDev-v0} & \textbf{HumanEval} \\
        \midrule
        Baseline           & 84.75 $\pm$ 1.94 & 84.25 $\pm$ 1.86 & 25.0 & 89.6 \\
        \hline
        Improved prompt    & 89.75 $\pm$ 1.44 & 89.00 $\pm$ 1.38 & 34.4 & 90.3 \\
        \hline
        New topology       & 85.50 $\pm$ 1.18 & 88.83 $\pm$ 1.51 & 40.6 & 91.5 \\
        \hline
    \end{tabular}
\end{table*}

% \melissa{DO NOT MENTION PROGRAMDEV HERE AS THIS IS NO LONGER THE PROGRAMDEV.}

\subsection{Effect of the interventions on \taxonomy}
\label{sec:interventions_effects_on_failure_modes_with_mast}

After carrying out the aforementioned interventions, we initially inspect the task completion rates as in Table \ref{tab:ag2-chatdev-accuracies}. However, \taxonomy offers us the opportunity to look beyond the task completion rates, and we can investigate the effects of these interventions on the failure mode distribution on these MASs (AG2 and ChatDev). As illustrated in Figures \ref{fig:intervention_comparison_threeway_ag2} and \ref{fig:intervention_comparison_threeway_chatdev}, we observe that both of these interventions cause a decrease across the different failure modes observed, and it is possible to conclude that topology-based changes are more effective than prompt-based changes for both systems. Moreover, this displays another usage of \taxonomy, which is as well as an analysis tool after execution, it can serve as a debugging tool for future improvements as it shows which failure modes particular augmentations to the system can solve or miss, guiding future intervention decisions.

\begin{figure}[h]
    \centering
    \includegraphics[width=0.9\linewidth]{figures/intervention_comparison_threeway_ag2_system_design.pdf}
    \caption{Effect of prompt and topology interventions on AG2 as captured by \taxonomy using the automated LLM-as-a-Judge}
    \label{fig:intervention_comparison_threeway_ag2}
\end{figure}

\begin{figure}[h]
    \centering
    \includegraphics[width=0.9\linewidth]{figures/intervention_comparison_threeway_chatdev_system_design.pdf}
    \caption{Effect of prompt and topology interventions on ChatDev as captured by \taxonomy using the automated LLM-as-a-Judge}
    \label{fig:intervention_comparison_threeway_chatdev}
\end{figure}




\newpage
\section{Analysis on Multi-Agent Systems with Open-Source Models}
\label{sec:analysis_on_opensource_models}

In this section, we also provide the analysis of failure modes on MetaGPT and ChatDev frameworks where the underlying LLMs are open-source models. In particular, we chose to use \verb|Qwen2.5-Coder-32B-Instruct| \cite{hui2024qwen2} and \verb|CodeLlama-7b-Instruct-hf| \cite{roziere2023code} models for these two frameworks. The analysis of failure modes is shown in Table \ref{tab:open_source_models}. This new analysis reveals two key findings:
\begin{itemize}
    \item There is a significant performance difference between the two open-source models. \verb|Qwen2.5-Coder-32B-Instruct| is substantially more robust than \verb|CodeLlama-7b-Instruct-hf| on these tasks, exhibiting far fewer failures overall.
    \item  Both open-source models show a higher frequency of failures compared to the leading closed-source models analyzed in our paper (GPT-4o and Claude-3). This suggests a performance gap and highlights important areas for future improvement in open-source models for multi-agent tasks.
\end{itemize}

\begin{table*}[ht!]
    \small
    \renewcommand{\arraystretch}{1.3}
    \setlength{\tabcolsep}{4pt}
    \centering
    \caption{Failure Mode Occurrences in 400 Traces with Open-Source Models. Results are grouped by model family (Qwen vs.\ CodeLlama) and development framework (ChatDev vs.\ MetaGPT).}
    \label{tab:failure-modes-open-source}
    \begin{tabular}{p{2.5cm} c c c c}
        \toprule
        \textbf{Failure Mode} & \multicolumn{2}{c}{\textbf{Qwen}} & \multicolumn{2}{c}{\textbf{CodeLlama}} \\
        \cmidrule(lr){2-3} \cmidrule(lr){4-5}
         & \textbf{ChatDev} & \textbf{MetaGPT} & \textbf{ChatDev} & \textbf{MetaGPT} \\
        \midrule
        1.1 & 35 & 12 & 76 & 94 \\
        \hline
        1.2 & 4 & 1 & 45 & 12 \\
        \hline
        1.3 & 96 & 35 & 97 & 99 \\
        \hline
        1.4 & 1 & 0 & 46 & 23 \\
        \hline
        1.5 & 94 & 3 & 97 & 76 \\
        \hline
        2.1 & 2 & 0 & 50 & 9 \\
        \hline
        2.2 & 1 & 4 & 16 & 15 \\
        \hline
        2.3 & 9 & 0 & 76 & 57 \\
        \hline
        2.4 & 0 & 0 & 2 & 0 \\
        \hline
        2.5 & 2 & 12 & 42 & 40 \\
        \hline
        2.6 & 20 & 16 & 93 & 18 \\
        \hline
        3.1 & 1 & 47 & 25 & 26 \\
        \hline
        3.2 & 16 & 51 & 67 & 55 \\
        \hline
        3.3 & 12 & 32 & 69 & 56 \\
        \hline
    \end{tabular}
\label{tab:open_source_models}
\end{table*}


\section{Correlation of Failure Modes with Different Statistics}
\label{sec:mast_correlations}

In this section, we provide how the failure modes in \taxonomy{} correlate with some important statistics, such as the actual task completion rates, and different benchmarks.


\subsection{How Indicative are Different Failure Modes of Actual Success?}

One important question to ask is, for traces where we know whether they succeeded or failed and if we do not provide the success or failure result to the LLM Annotator, how do different failure modes correlate with actual task success and failures? In particular, we want to measure how indicative of the failure modes given by the LLM Annotator on actual task completions. The analysis on ChatDev and MetaGPT are provided in Table \ref{tab:failure_mode_occurrence_on_successful_traces}.  This new analysis reveals three key findings:
\begin{itemize}
    \item Successful runs are not failure-free. Our results show that failures occur in both successful and failed traces, but failed traces have a higher overall frequency of failures. This confirms that a higher number of failures do signal a higher chance of final task failure.
    \item Some failures are more "fatal" than others. The data shows a clear distinction between failure types. Certain failures, such as 1.5 Unaware of Termination Conditions and 2.4 Information Withholding, appear almost exclusively in failed runs, suggesting they are critical bugs that are highly likely to derail the task.
    \item Some failures are non-fatal. In contrast, verification-related failures like 3.2 No or Incomplete Verification and 3.3 Incorrect Verification appear frequently even in successful runs. This suggests that while these systems can complete some tasks, their verification process still contains flaws. MAST identifies such systemic weaknesses, even when they do not cause an immediate task failure.
\end{itemize}


\begin{table*}[ht!]
    \small
    \renewcommand{\arraystretch}{1.3}
    \setlength{\tabcolsep}{3.5pt}
    \centering
    \caption{Failure mode occurrence rates for ChatDev and MetaGPT on successful and unsuccessful examples.}
    \label{tab:per-failure-mode-outcomes}
    \begin{tabular}{p{3.0cm} *{14}{c}}
        \toprule
        \textbf{MAS Framework} & \multicolumn{14}{c}{\textbf{Failure mode occurrence (\%)}} \\
        \cmidrule(lr){2-15}
         & \textbf{1.1} & \textbf{1.2} & \textbf{1.3} & \textbf{1.4} & \textbf{1.5} & \textbf{2.1} & \textbf{2.2} & \textbf{2.3} & \textbf{2.4} & \textbf{2.5} & \textbf{2.6} & \textbf{3.1} & \textbf{3.2} & \textbf{3.3} \\
        \midrule
        ChatDev Success & 20.0 & 0.0 & 20.0 & 0.0 & 0.0 & 0.0 & 10.0 & 10.0 & 0.0 & 0.0 & 10.0 & 0.0 & 10.0 & 20.0 \\
        \hline
        ChatDev Fail    & 25.0 & 0.0 & 20.0 & 5.0 & 10.0 & 5.0 & 5.0  & 5.0  & 5.0 & 0.0 & 15.0 & 5.0 & 10.0 & 25.0 \\
        \hline
        MetaGPT Success & 33.3 & 0.0 & 16.7 & 0.0 & 0.0 & 0.0 & 8.3  & 8.3  & 0.0 & 0.0 & 8.3  & 0.0 & 16.7 & 16.7 \\
        \hline
        MetaGPT Fail    & 16.7 & 0.0 & 22.2 & 5.6 & 11.1 & 5.6 & 5.6  & 5.6  & 5.6 & 0.0 & 16.7 & 5.6 & 5.6  & 27.8 \\
        \hline
    \end{tabular}
        \label{tab:failure_mode_occurrence_on_successful_traces}

\end{table*}

\subsection{Failure Mode Occurrence Rates on Different Benchmarks}

We also analyze the rate of failure modes on different benchmarks, ranging from general question answering like MMLU to math problems like GSM and harder reasoning problems in OlympiadBench. For this, we fixed the MAS framework (AG2) and the model (GPT-4o) while varying the benchmark. The results are presented in Table \ref{tab:failure_mode_rates_on_benchmarks}. The failure rate is normalized by the number of traces.
We observed that the more challenging the benchmark is (e.g., Olympiad vs. GSM), the higher the failure rate. The distribution also changes significantly. While the failure profiles for MMLU and Olympiad are similar, the GSM benchmark results in a much lower number of Inter-Agent Misalignment and Specification failures.

\begin{table*}[ht!]
    \small
    \renewcommand{\arraystretch}{1.3}
    \setlength{\tabcolsep}{4pt}
    \centering
    \caption{Failure Category Rates on Different Benchmarks}
    \label{tab:fc-components}
    \begin{tabular}{p{2cm} p{1.4cm} c c c}
        \toprule
        \textbf{MAS / LLM} & \textbf{Benchmark} & \textbf{FC1: System Design} & \textbf{FC2: Inter-Agent Misalignment} & \textbf{FC3: Verification} \\
        \midrule
        AG2 / GPT-4o & GSM       & 0.53 & 1.33 & 0.37 \\
        AG2 / GPT-4o & MMLU      & 1.06 & 1.01 & 0.60 \\
        AG2 / GPT-4o & Olympiad  & 1.19 & 1.21 & 0.67 \\
        \bottomrule
    \end{tabular}
    \label{tab:failure_mode_rates_on_benchmarks}
\end{table*}




\section{LLM Annotator Cost}
\label{sec:llm_judge_cost}

We have analyzed the API costs for our LLM-as-a-Judge pipeline across all traces in our study. The average cost across all MAS frameworks is $\$1.8$, and the average cost per MAS highly depends on the length of the traces. The cost breakdown by MAS framework (normalized by the number of traces collected) is shown in Table \ref{tab:llm_judge_cost_breakdown}.

\begin{table*}[ht!]
    \small
    \renewcommand{\arraystretch}{1.25}
    \setlength{\tabcolsep}{6pt}
    \centering
    \caption{Average failure cost by MAS framework.}
    \label{tab:avg-cost-by-mas}
    \begin{tabular}{p{3.0cm} r}
        \toprule
        \textbf{MAS} & \textbf{Average Cost} \\
        \midrule
        AppWorld    & 0.3740 \\
        HyperAgent  & 0.9695 \\
        AG2         & 1.1656 \\
        ChatDev     & 2.1272 \\
        MetaGPT     & 2.4455 \\
        MagenticOne & 1.3056 \\
        OpenManus   & 4.1409 \\
        \bottomrule
    \end{tabular}
    \label{tab:llm_judge_cost_breakdown}
\end{table*}





\newpage
\section{AG2 - MathChat Scenario}
\label{appendix:ag2}

\subsection{Initial prompt}
\label{sec:ag2-mathchat-original-prompt}
\begin{lstlisting}[
    basicstyle=\footnotesize\ttfamily,
    columns=flexible,
    breaklines=true
]
Let's use Python to solve a math problem.

Query requirements:
You should always use the 'print' function for the output and use fractions/radical forms instead of decimals.
You can use packages like sympy to help you.
You must follow the formats below to write your code:

```python
# your code
```

First state the key idea to solve the problem. You may choose from three ways to solve the problem:
Case 1: If the problem can be solved with Python code directly, please write a program to solve it. You can enumerate all possible arrangements if needed.
Case 2: If the problem is mostly reasoning, you can solve it by yourself directly.
Case 3: If the problem cannot be handled in the above two ways, please follow this process:
1. Solve the problem step by step (do not over-divide the steps).
2. Take out any queries that can be asked through Python (for example, any calculations or equations that can be calculated).
3. Wait for me to give the results.
4. Continue if you think the result is correct. If the result is invalid or unexpected, please correct your query or reasoning.

After all the queries are run and you get the answer, put the answer in \\boxed{}.

Problem:
\end{lstlisting}


\subsection{Structured prompt with verification section}
\label{sec:ag2-mathchat-improved-prompt}
\begin{lstlisting}[
    basicstyle=\footnotesize\ttfamily,
    columns=flexible,
    breaklines=true
]
Let's use Python to tackle a math problem effectively.

Query Requirements:
1. Output Format: Always utilize the print function for displaying results. Use fractions or radical forms instead of decimal numbers.
2. Libraries: You are encouraged to use packages such as sympy to facilitate calculations.

Code Formatting:
Please adhere to the following format when writing your code:
```python
# your code
```

Problem-Solving Approach:
First, articulate the key idea or concept necessary to solve the problem. You can choose from the following three approaches:
Case 1: Direct Python Solution. If the problem can be solved directly using Python code, write a program to solve it. Feel free to enumerate all possible arrangements if necessary.
Case 2: Reasoning-Based Solution. If the problem primarily involves reasoning, solve it directly without coding.
Case 3: Step-by-Step Process. If the problem cannot be addressed using the above methods, follow this structured approach:
1. Break down the problem into manageable steps (avoid excessive granularity).
2. Identify any queries that can be computed using Python (e.g., calculations or equations).
3. Await my input for any results obtained.
4. If the results are valid and expected, proceed with your solution. If not, revise your query or reasoning accordingly.

Handling Missing Data:
If a problem is deemed unsolvable due to missing data, return \boxed{'None'}. 
Ensure that only numerical values are placed inside the \boxed{}; any accompanying words should be outside.

Verification Steps:
Before presenting your final answer, please complete the following steps:
1. Take a moment to breathe deeply and ensure clarity of thought.
2. Verify your solution step by step, documenting each part of the verification process in a designated VERIFICATION section.
3. Once you are confident in your verification and certain of your answer, present your final result in the format \boxed{_you_answer_}, ensuring only numbers are inside.

Problem Statement: 
\end{lstlisting}


\subsection{Agent Problem Solver's System Prompt}
\label{sec:ag2-mathchat-new-topology-problem-solver}
\begin{lstlisting}[
    basicstyle=\footnotesize\ttfamily,
    columns=flexible,
    breaklines=true
]
You are Agent Problem Solver, and your role is to collaborate with other agents to address various challenges. 

For each problem, please follow these steps:
1. **Document Your Solution**: Write your solution step by step, ensuring it is independent of the solutions provided by other agents.
2. **Engage in Discussion**: Once you have outlined your solution, discuss your approach and findings with the other agents.
\end{lstlisting}


\subsection{Agent Coder's System Prompt}
\label{sec:ag2-mathchat-new-topology-coder}
\begin{lstlisting}[
    basicstyle=\footnotesize\ttfamily,
    columns=flexible,
    breaklines=true
]
You are Agent Code Executor. You can solve problems only writing commented Python code.

For each problem, please follow these steps:
1. **Develop Your Solution**: Write your solution in Python code, detailing each step independently from the solutions provided by other agents.
2. **Utilize SymPy**: Feel free to use the SymPy package to facilitate calculations and enhance your code's efficiency.
3. **Display Results**: Ensure that you **print the final result at the end of your Python code** (e.g., `print(_result_)`).
4. **Engage in Discussion**: After obtaining the result from your Python code, discuss your findings with the other agents.

Always format your Python code within: 
```python
# your code here
print(_result_)
``` 

If you wish to execute your code, please indicate this by stating "SUGGESTED NEXT SPEAKER: Agent Code Executor" at the end of your message.
\end{lstlisting}


\subsection{Agent Verifier's System Prompt}
\label{sec:ag2-mathchat-new-topology-verifier}
\begin{lstlisting}[
    basicstyle=\footnotesize\ttfamily,
    columns=flexible,
    breaklines=true
]
You are Agent Verifier.    

Your role is to critically evaluate the solutions proposed by other agents step by step and provide a final solution. 
1. **Solution Requirement**: Before making any decisions, ensure you have received solutions from both Agent Code Executor and Agent Problem Solver. If either proposed solution is missing, do not draw any conclusions; instead, suggest the next speaker by stating: SUGGESTED NEXT SPEAKER: _suggested_agent_name_.
2. **Avoid Assumptions**: Pay attention to the variables provided in the original problem statement versus those assumed by the agents. **Assumed values are not valid for the solution** and can lead to inaccuracies. Never base your solution on assumed values. Always base your solution on the explicitly given variables to ensure correctness. If a problem is deemed unsolvable due to missing information, return: **SOLUTION_FOUND \\boxed{'None'}**.
3. **Evaluating Conflicting Solutions**: If different answers are presented during the discussion, choose the most appropriate solution based on your evidence or initiate further discussion to clarify.
4. **Final Solution Declaration**: When you are confident about the final solution, return it as follows: **SOLUTION_FOUND \\boxed{_solution_value_here_}**. Ensure that only numerical values are placed inside the \\boxed{}; any accompanying text should be outside.
\end{lstlisting}


\subsection{AG2 - MathChat: an example of conversation with three specialized agents}
\label{sec:ag2-mathchat-new-topology-example}
\begin{lstlisting}[
    basicstyle=\footnotesize\ttfamily,
    columns=flexible,
    breaklines=true
]
[Agent Verifier]
Sara wants to buy herself a new jacket and 2 pairs of shoes. The jacket she wants costs $120 and each pair of shoes cost $100. Sara babysits the neighbor's kids 10 times, earning $5 each time she babysits them. Her parents pay her $10 each time she mows the lawn. If Sara already had $100 saved before she started babysitting, how many times must she mow the lawn before she can afford the jacket and shoes?

[Agent Problem Solver]
First, let's break down the problem:,
- The costs are: $120 for the jacket, $100 for each pair of shoes, and she wants 2 pairs. So the total cost is $120 + 2*$100 = $320,
- Sara earns $5 each time she babysits, and she has done this 10 times, so she has earned $5*10 = $50 from babysitting.
- Sara also has $100 saved from before.
- The amount Sara has after babysitting and her savings is $100 + $50 = $150.
- The amount Sara still needs is the total cost minus what she has, which is $320 - $150 = $170.
- Sara earns $10 for every time she mows the lawn, to determine how many times she must do this, we divide the amount she needs by the amount she earns each time: $170 / $10 = 17 times.
So, Sara must mow the lawn 17 times to afford the jacket and shoes.

[Agent Verifier]
SUGGESTED NEXT SPEAKER: Agent_Code_Executor

[Agent Coder]
Sure, I'll write a Python code for the calculation.
```python
# cost of items
jacket_cost = 120
shoes_cost = 100 * 2  # since she wants 2 pairs of shoes
# total cost of items
total_cost = jacket_cost + shoes_cost
# earnings from babysitting
babysitting_earning = 5 * 10  # $5 each time, done 10 times
# initial money
initial_money = 100
# total money after babysitting
total_money_after_babysitting = initial_money + babysitting_earning
# amount still needed
amount_needed = total_cost - total_money_after_babysitting
# how much sara earns from mowing each time
mowing_earning = 10
# calculate how many times she needs to mow the lawn
num_times_mowing = amount_needed // mowing_earning  # using integer division for whole number of times
# print result
print(num_times_mowing)
```
When you execute this, the result should be 17, confirming Agent Problem Solver's solution.

[Agent Verifier]
exitcode: 0 (execution succeeded)
Code output: 17

[Agent Problem Solver]
Agent Code Executor's Python program has provided the solution - Sara must mow the lawn 17 times to afford the jacket and shoes, which supports Agent Problem Solver's proposal. Both paths lead us to the same conclusion, which approves the accuracy of our solution. 
Great team work, Agents.

[Agent Verifier]
SOLUTION_FOUND \\boxed{17}
\end{lstlisting}










\newpage
\section{ChatDev Prompt Interventions}
\label{appendix:chatdev}


\subsection{Initial Prompt}
\label{appendix:chatdev_inits}

\subsubsection{Programmer}
\begin{lstlisting}[
    basicstyle=\footnotesize\ttfamily,
    columns=flexible,
    breaklines=true
]
"{chatdev_prompt}",
    "You are Programmer. we are both working at ChatDev. We share a common interest in collaborating to successfully complete a task assigned by a new customer.",
    "You can write/create computer software or applications by providing a specific programming language to the computer. You have extensive computing and coding experience in many varieties of programming languages and platforms, such as Python, Java, C, C++, HTML, CSS, JavaScript, XML, SQL, PHP, etc,.",
    "Here is a new customer's task: {task}.",
    "To complete the task, you must write a response that appropriately solves the requested instruction based on your expertise and customer's needs."
 
\end{lstlisting}
\subsubsection{Code Reviewer}
\begin{lstlisting}[
    basicstyle=\footnotesize\ttfamily,
    columns=flexible,
    breaklines=true
]
"{chatdev_prompt}",
    "You are Code Reviewer. we are both working at ChatDev. We share a common interest in collaborating to successfully complete a task assigned by a new customer.",
    "You can help programmers to assess source codes for software troubleshooting, fix bugs to increase code quality and robustness, and offer proposals to improve the source codes.",
    "Here is a new customer's task: {task}.",
    "To complete the task, you must write a response that appropriately solves the requested instruction based on your expertise and customer's needs."

\end{lstlisting}
\subsubsection{Software Test Engineer}
\begin{lstlisting}[
    basicstyle=\footnotesize\ttfamily,
    columns=flexible,
    breaklines=true
]
  "{chatdev_prompt}",
    "You are Software Test Engineer. we are both working at ChatDev. We share a common interest in collaborating to successfully complete a task assigned by a new customer.",
    "You can use the software as intended to analyze its functional properties, design manual and automated test procedures to evaluate each software product, build and implement software evaluation test programs, and run test programs to ensure that testing protocols evaluate the software correctly.",
    "Here is a new customer's task: {task}.",
    "To complete the task, you must write a response that appropriately solves the requested instruction based on your expertise and customer's needs."
 
\end{lstlisting}
\subsubsection{Chief Executive Officer}
\begin{lstlisting}[
    basicstyle=\footnotesize\ttfamily,
    columns=flexible,
    breaklines=true
]
"{chatdev_prompt}",
    "You are Chief Executive Officer. Now, we are both working at ChatDev and we share a common interest in collaborating to successfully complete a task assigned by a new customer.",
    "Your main responsibilities include being an active decision-maker on users' demands and other key policy issues, leader, manager, and executor. Your decision-making role involves high-level decisions about policy and strategy; and your communicator role can involve speaking to the organization's management and employees.",
    "Here is a new customer's task: {task}.",
    "To complete the task, I will give you one or more instructions, and you must help me to write a specific solution that appropriately solves the requested instruction based on your expertise and my needs."
    
\end{lstlisting}
\subsubsection{Chief Technology Officer}
\begin{lstlisting}[
    basicstyle=\footnotesize\ttfamily,
    columns=flexible,
    breaklines=true
]
 "{chatdev_prompt}",
    "You are Chief Technology Officer. we are both working at ChatDev. We share a common interest in collaborating to successfully complete a task assigned by a new customer.",
    "You are very familiar to information technology. You will make high-level decisions for the overarching technology infrastructure that closely align with the organization's goals, while you work alongside the organization's information technology (\"IT\") staff members to perform everyday operations.",
    "Here is a new customer's task: {task}.",
    "To complete the task, You must write a response that appropriately solves the requested instruction based on your expertise and customer's needs."
 
\end{lstlisting}
\subsection{Modified System Prompts}

\subsubsection{Programmer}
\begin{lstlisting}[
    basicstyle=\footnotesize\ttfamily,
    columns=flexible,
    breaklines=true
]
 "{chatdev_prompt}",
      "You are a Programmer at ChatDev. Your primary responsibility is to develop software applications by writing code in various programming languages. You have extensive experience in languages such as Python, Java, C++, JavaScript, and others. You translate project requirements into functional and efficient code.",
      "You report to the technical lead or CTO and collaborate with other programmers and team members.",
      "Here is a new customer's task: {task}.",
      "To complete the task, you will write code to implement the required functionality, ensuring it meets the customer's specifications and quality standards."

\end{lstlisting}
\subsubsection{Software Test Engineer}
\begin{lstlisting}[
    basicstyle=\footnotesize\ttfamily,
    columns=flexible,
    breaklines=true
]
 "{chatdev_prompt}",
      "You are a Software Test Engineer at ChatDev. Your primary responsibility is to design and execute tests to ensure the quality and functionality of software products. You develop test plans, create test cases, and report on software performance. You identify defects and collaborate with the development team to resolve them.",
      "You need to ensure that the software is working as expected and meets the customer's requirements.",
      "Check the edge cases and special cases and instances for the task we are doing. Do not miss any cases. Do not suffice with generic and superficial cases.",
      "You report to the technical lead or CTO and collaborate with programmers and code reviewers.",
      "Here is a new customer's task: {task}.",
      "To complete the task, you will design and implement test procedures, report issues found, and verify that the software meets the customer's requirements."

\end{lstlisting}
\subsubsection{Code Reviewer}
\begin{lstlisting}[
    basicstyle=\footnotesize\ttfamily,
    columns=flexible,
    breaklines=true
]
"{chatdev_prompt}",
      "You are a Code Reviewer at ChatDev. Your primary responsibility is to review and assess source code written by programmers. You ensure code quality by identifying bugs, optimizing performance, and enforcing coding standards. You provide constructive feedback to improve software robustness.",
      "You report to the technical lead or CTO and work closely with programmers.",
      "Here is a new customer's task: {task}.",
      "To complete the task, you will review the code submitted by programmers, identify issues, and suggest improvements to meet quality standards."

\end{lstlisting}
\subsubsection{Chief Executive Officer}
\begin{lstlisting}[
    basicstyle=\footnotesize\ttfamily,
    columns=flexible,
    breaklines=true
]
"{chatdev_prompt}",
      "You are the Chief Executive Officer (CEO) of ChatDev. Your primary responsibilities include making high-level decisions about policy and strategy, overseeing the overall operations and resources of ChatDev, and acting as the main point of communication between the board and corporate operations.",
      "As the CEO, you have the authority to make final decisions and terminate conversations when appropriate.",
      "Here is a new customer's task: {task}.",
      "To complete the task, you will provide strategic guidance and instructions to your team, ensuring that the solution meets the customer's needs and aligns with the company's objectives."

\end{lstlisting}
\subsubsection{Chief Technology Officer}
\begin{lstlisting}[
    basicstyle=\footnotesize\ttfamily,
    columns=flexible,
    breaklines=true
]
 "{chatdev_prompt}",
      "You are the Chief Technology Officer (CTO) of ChatDev. Your primary responsibilities include overseeing all technical aspects of the company. You establish the company's technical vision and lead technological development, ensuring that technology resources align with the company's business needs.",
      "You report to the CEO and collaborate with other executives to integrate technology into the company's strategy.",
      "Here is a new customer's task: {task}.",
      "To complete the task, you will develop the technical strategy and guide your team to ensure the solution meets the customer's needs and adheres to technological standards."

\end{lstlisting}


% \newpage
% \section{ProgramDev}
% \melissa{here we also release ProgramDev task for full transparency. ProgramDev is intended to be a task that university level CS student should be able to implement with open access to internet. However, the model struggle to implement. etc...}

\newpage
\section{Examples of Different Failure Modes}
\label{appendix:examples}

\subsection{Example for FM-1.3 Step repetition, FM-2.1 Conversation reset, FM-1.2 Disobey role specification}\label{ex:fm2.2}

\begin{itemize}
    \item \textbf{MAS:} Hyperagent
    \item \textbf{Task:} Matplotlib visualization code bug fix.
    \item \textbf{What happened:} The Navigator internally proposes a solution without communicating it to the Planner. They did not manage to replicate the issue described and started to unrelated steps to find something to edit. At the end the Planner decided to 'Wait for the user's response'. 

    \item \textbf{Where it went wrong:}
    \begin{lstlisting}[
        basicstyle=\footnotesize\ttfamily,
        columns=flexible,
        breaklines=true
    ]
[...]
HyperAgent_matplotlib__matplotlib-23563 - INFO - Planner's Response: Thought: To address this issue, we need to understand the root cause of the 'Line3D' object not having the '_verts3d' attribute. This error typically occurs when the data passed to the plot function is invalid or when there's an issue with the 3D rendering. Let's start by examining the code and reproducing the issue.
[...]
HyperAgent_matplotlib__matplotlib-23563 - INFO - Planner's Response: Thought: To address this issue, we need to understand the root cause of the 'Line3D' object not having the '_verts3d' attribute. This error typically occurs when the data passed to the plot function is invalid or when there's an issue with the 3D rendering. Let's start by examining the code and reproducing the issue.
[...]
Intern Name: Codebase Editor
Subgoal: Please prepare a response for the user with the following content:
1. Explain the identified issue and recommend updating matplotlib to version 3.5.1 or later.
2. Provide a modified version of their code that attempts to work around the issue.
3. Suggest an alternative plotting method using scatter plots instead of lines.
[...]
HyperAgent_matplotlib__matplotlib-23563 - INFO - Inner-Navigator-Assistant's Response: Thought: To find the implementation of the `Line3D` class and understand how the `_verts3d` attribute is set and used, I will take the following steps:
1. Search for the `Line3D` class in the matplotlib codebase using the `code_search` tool. [...]
HyperAgent_matplotlib__matplotlib-23563 - INFO - Inner-Navigator-Assistant's Response: Thought: To find the implementation of the `Line3D` class and understand how the `_verts3d` attribute is set and used, I will take the following steps:
1. Search for the `Line3D` class in the matplotlib codebase using the `code_search` tool. 
[...]
    \end{lstlisting}
\end{itemize}







\subsection{Example for FM-1.5 Unaware of stopping conditions}\label{ex:fm1.5}
\begin{itemize}
    \item \textbf{MAS: AG2} 
    \item \textbf{Task: Solving a mathematical problem} 
    \item \textbf{What happened}: The solution is correct. The mathproxyagent keeps asking to continue even if it was clear that the problem was unsolvable. 
    \item \textbf{Where it went wrong:}
    \begin{lstlisting}[
        basicstyle=\footnotesize\ttfamily,
        columns=flexible,
        breaklines=true
    ]
[...]
[Assistant] This problem is a reasoning problem. The information provided in the problem statement does not give any specifics about the amount of ribbon Monica has or the size of the gift bows. Therefore, this problem cannot be solved with the information provided. More context or information is needed, such as \"Monica has 120 inches of ribbon\" or \"Each bow uses X inches of ribbon\"."
[mathproxyagent] Continue. Please keep solving the problem until you need to query. (If you get to the answer, put it in \\boxed{}.)
[Assistant] In the current problem, I don't have enough information to solve the problem. [...]
[mathproxyagent] Continue. Please keep solving the problem until you need to query. (If you get to the answer, put it in \\boxed{}.)
[Assistant] I'm sorry, but I really need more information to solve this problem. 
[mathproxyagent] Continue. Please keep solving the problem until you need to query. (If you get to the answer, put it in \\boxed{}.)
[...]
    \end{lstlisting}
\end{itemize}



\subsection{Example for FM-2.3 Task derailment}\label{ex:fm2.5}
\begin{itemize}
    \item \textbf{MAS:} AG2 
    \item \textbf{Task:} Solving a mathematical problem.
    \item \textbf{What happened:} The answer was given the the question. The agents solved another problem. 
    \item \textbf{Where it went wrong:}
    \begin{lstlisting}[
        basicstyle=\footnotesize\ttfamily,
        columns=flexible,
        breaklines=true
    ]
[mathproxyagent] Let's use Python to solve a math problem. [...] Courtney attended a concert and reported that the audience was 48 in number. However, Kelly went to the same concert and said that Courtney had made the mistake of overstating the number of people in attendance by 20%. If Kelly was right, 40 people really attended the concert. What was the audience size according to Courtney?"
[assistant] First, the key idea here is calculating [...] 
```python
import sympy as sp
C = 48  # Courtney's estimation
R = sp.Symbol('R')  # real audience number
# equation using the information given by Kelly
equation = sp.Eq(C, R + 0.2 * R)
real_audience = sp.solve(equation, R)[0]
print(f\"The real audience was: {real_audience}\")
```
We subtract the estimated audience size from the real audience size to find the number of people Courtney overstated.
```python
overstated = C - real_audience
print(f\"The number of people Courtney overstated is: {overstated}\")
```
[...]
[assistant] [...] So, the final answer is \\boxed{8}.
    \end{lstlisting}
\end{itemize}



\subsection{Example for FM-2.2 Fail to ask for clarification}\label{ex:fm1.4}
\begin{itemize}
    \item \textbf{MAS:} AppWorld 
    \item \textbf{Task:} Give a 5-star rating to all songs in my Spotify playlists. 
    \item \textbf{What happened:} The agent used the password as token. 
    \item \textbf{Where it went wrong:}
    \begin{lstlisting}[
        basicstyle=\footnotesize\ttfamily,
        columns=flexible,
        breaklines=true
    ]
[...]
[Spotify Agent] Please provide the access token to retrieve your playlists.
[Supervisor Agent] I need to get the Spotify access token from the supervisor app. Let's retrieve the account credentials to find the access token. print(apis.supervisor.show_account_passwords())
[...]
    \end{lstlisting}
\end{itemize}



\subsection{Example for FM-3.2 No or incomplete verification, }\label{ex:fm4.2}
\begin{itemize}
    \item \textbf{MAS:} AG2
    \item \textbf{Task:} Mathematical problem solving
    \item \textbf{What happened:} Confused the number of fish with the cost of them.
    \item \textbf{Where it went wrong:}
    \begin{lstlisting}[
        basicstyle=\footnotesize\ttfamily,
        columns=flexible,
        breaklines=true
    ]
[mathproxyagent] [...] Scarlett found an aquarium for $10.00 at a yard sale.  At the pet store, she bought 2 bags of rocks for $2.50 each and 3 pieces of coral at $2.00 apiece.  She bought 20 fish and she needed fish food that cost $2.00.  How much did she spend?
[assistant] [... ] fish_cost = 20.00 [...] total_cost = aquarium_cost + 2*rock_cost + 3*coral_cost + fish_cost + fish_food_cost [...] print(total_cost)
[mathproxyagent] 43.0 
[assistant] So, Scarlett spent a total of $43.00. So, the answer is \\boxed{43}.
    \end{lstlisting}
\end{itemize}



\subsection{Example for FM-2.3 Task derailment, FM-3.2 No or incomplete verification.}\label{ex:fm2.5-2}\label{ex:fm4.2-2}
\begin{itemize}
    \item \textbf{MAS:} AG2
    \item \textbf{Task:} Solving a mathematical problem. 
    \item \textbf{What happened:} Found the solution correctly, then it decides to answer to another problem. 
    \item \textbf{Where it went wrong:}
    \begin{lstlisting}[
        basicstyle=\footnotesize\ttfamily,
        columns=flexible,
        breaklines=true
    ]
[mathproxyagent] [...] John goes to the market with $10,000 in his pocket. He buys 3 goats for $500 each and 2 cows for $1500 each. He also considers buying a horse for $2000 but decides against it. How much money did he spend at the market?
[...]
[assistant] Great, the total spent by John on buying goats and cows is $4500. To find the remaining money John had after his purchases, we subtract the total spent from the initial amount of money he had. [...] This code will output how much money John had left after he made his purchases at the market.
[...]
[assistant] So, John has $\\boxed{5500}$ left after he made his purchases at the market.
    \end{lstlisting}
\end{itemize}



\subsection{Example for FM-3.2 No or incomplete verification}\label{ex:fm1.3}\label{ex:fm4.3}
\begin{itemize}
    \item \textbf{MAS:} AG2
    \item \textbf{Task:} Solving a mathematical problem. 
    \item \textbf{What happened:} They did not verify the answer. They counted one number more even if the numbers were listed. They did not verify the result even if it was straightforward.  
    \item \textbf{Where it went wrong:}
    \begin{lstlisting}[
        basicstyle=\footnotesize\ttfamily,
        columns=flexible,
        breaklines=true
    ]
[mathproxyagent] [...] Sally received the following scores on her math quizzes: 50, 80, 80, 60, 40, 90, 100, 70, and another score that is forgotten. If the mean score is 70, what's the forgortten score?
[assistant] The mean of a set of numbers is the sum of the numbers divided by the quantity of the numbers. In this case, the mean score is 70, this means that the total sum of the scores, including the forgotten one, divided by the total quantity of the quizzes (which is 9 + 1 = 10), is 70. [...] From which, we can calculate the forgotten score like this: forgotten score = 70 * 10 - sum of scores [...]
[mathproxyagent] Continue. Please keep solving the problem until you need to query. (If you get to the answer, put it in \\boxed{}.)
[...]
[assistant] Therefore, the forgotten score is \\boxed{130}. 
    \end{lstlisting}
\end{itemize}



\subsection{Example for FM-1.2 Disobey role specification, FM-2.4 Information withholding}\label{ex:fm2.6}
\begin{itemize}
    \item \textbf{MAS:} HyperAgent 
    \item \textbf{Task:} Solving a bug of matplotlib.
    \item \textbf{What happened:} The Navigator thinks also about possible solutions (not required) but then it does not communicate it to the Planner. At certain point the Executor states 'The above output is just an example' but the Planner thinks that the test are passed 'Thank you for running the test and providing the results'. 
    \item \textbf{Where it went wrong:}
    \begin{lstlisting}[
        basicstyle=\footnotesize\ttfamily,
        columns=flexible,
        breaklines=true
    ]
[...]
HyperAgent_matplotlib__matplotlib-25433 - INFO - Inner-Navigator-Assistant's Response: Thought: [...] To work around this issue, a possible solution is to update [...] Here's an example of how you can modify the code to achieve this: [...] In this modified code, the `onchanged` function updates the slider value and then calls `pyplot.draw()` to redraw the figure. This should prevent the input to the widgets from being blocked. 
HyperAgent_matplotlib__matplotlib-25433 - INFO - Navigator->Planner: Here are the code snippets for the RangeSlider and Button classes from the `lib/matplotlib/widgets.py` file: [...]
    \end{lstlisting}
\end{itemize}


\subsection{Example for FM-1.1 Disobey task specification, FM-3.3 Incorrect verification}\label{ex:fm1.1}
\begin{itemize}
    \item \textbf{MAS:} MetaGPT
    \item \textbf{Task:} Implementing chess. 
    \item \textbf{What happened:} Missing coverage of movement restrictions for non-pawn pieces in initial tests, and lack of verification for complete chess rules.
    \item \textbf{Where it went wrong:}
    \begin{lstlisting}[
        basicstyle=\footnotesize\ttfamily,
        columns=flexible,
        breaklines=true
    ]
2025-01-12 14:46:45.578 | INFO     | metagpt.roles.role:_act:391 - Alice(SimpleCoder): to do SimpleWriteCode(SimpleWriteCode)
```python
def main():
    import itertools
    class ChessPiece:
        def __init__(self, color, name):
            self.color = color
            self.name = name
        def __str__(self):
            return f"{self.color[0]}{self.name[0]}"
        def valid_moves(self, position, board):
            moves = []
            if self.name == "Pawn":
                direction = 1 if self.color == "White" else -1
                forward = (position[0] + direction, position[1])
                if 0 <= forward[0] <= 7 and board[forward[0]][forward[1]] is None:
                    moves.append(forward)
                if position[0] == (6 if self.color == "White" else 1):  # Starting position
                    double_forward = (position[0] + 2 * direction, position[1])
                    if board[double_forward[0]][double_forward[1]] is None:
                        moves.append(double_forward)
            # Add more piece logic here
            return moves
[...]
def test_knight_valid_moves(setup_board):
    board = setup_board
    # Place a knight on the board
    board.board[0][1] = ChessPiece("White", "Knight")
    # Test if knight can move to a valid position
    assert board.move_piece(get_position('b1'), get_position('c3')) == True
def test_bishop_valid_moves(setup_board):
    board = setup_board
    # Place a bishop on the board
    board.board[0][2] = ChessPiece("White", "Bishop")
    # Test if bishop can move diagonally
    assert board.move_piece(get_position('c1'), get_position('a3')) == True
def test_rook_valid_moves(setup_board):
    board = setup_board
    # Place a rook on the board
    board.board[0][0] = ChessPiece("White", "Rook")
    # Test if rook can move vertically
    assert board.move_piece(get_position('a1'), get_position('a4')) == True
def test_queen_valid_moves(setup_board):
    board = setup_board
    # Place a queen on the board
    board.board[0][3] = ChessPiece("White", "Queen")
    # Test if queen can move diagonally
    assert board.move_piece(get_position('d1'), get_position('f3')) == True
def test_king_valid_moves(setup_board):
    board = setup_board
    # Place a king on the board
    board.board[0][4] = ChessPiece("White", "King")
    # Test if king can move one square in any direction
    assert board.move_piece(get_position('e1'), get_position('e2')) == True
[...]
    \end{lstlisting}
\end{itemize}



\subsection{Example for FM-2.6 Reasoning-action mismatch}\label{ex:fm1.2}
\begin{itemize}
    \item \textbf{MAS:} HyperAgent 
    \item \textbf{Task:} Solving a bug of pylint. 
    \item \textbf{What happened:} Initially, the navigator has already found the solution but it tells other things to planner. 
    \item \textbf{Where it went wrong:} 
    \begin{lstlisting}[
        basicstyle=\footnotesize\ttfamily,
        columns=flexible,
        breaklines=true
    ]
[...]
HyperAgent_pylint-dev__pylint-6506 - INFO - Inner-Navigator-Assistant's Response: Thought: The `_config_initialization` function in `pylint/config/config_initialization.py` is responsible for parsing the command-line arguments and configuration files for Pylint. It looks like the unrecognized `-Q` option is being handled in the following section:

```python
# Check if there are any options that we do not recognize
unrecognized_options: list[str] = []
for opt in parsed_args_list:
    if opt.startswith(\"--\"):
        unrecognized_options.append(opt[2:])
    elif opt.startswith(\"-\"):
        unrecognized_options.append(opt[1:])
if unrecognized_options:
    msg = \", \".join(unrecognized_options)
    linter.add_message(\"unrecognized-option\", line=0, args=msg)
    raise _UnrecognizedOptionError(options=unrecognized_options)
```

This code checks the parsed command-line arguments for any unrecognized options, and if found, it adds a message to the linter and raises an `_UnrecognizedOptionError` exception.

The issue mentioned in the GitHub report is that the traceback is being printed, which is not the expected behavior. The expected behavior should be to print a more user-friendly error message, similar to what the top part of the output already provides.

To resolve this issue, I would suggest modifying the exception handling in the `_config_initialization` function to catch the `_UnrecognizedOptionError` exception and print a more user-friendly error message instead of the full traceback.

Final Answer:
The root cause of the issue is in the `pylint/config/config_initialization.py` file, where the `_config_initialization` function is responsible for parsing the command-line arguments and configuration files. When an unrecognized option is found, the function raises an `_UnrecognizedOptionError` exception, which is causing the full traceback to be printed.

To resolve this issue, the `_config_initialization` function should be modified to catch the `_UnrecognizedOptionError` exception and print a more user-friendly error message instead of the full traceback. Here's an example of how the code could be updated:

```python
try:
    linter._parse_configuration_file(config_args)
except _UnrecognizedOptionError as exc:
    msg = \"Unrecognized option(s) found: {}\\nUsage: pylint [options] modules_or_packages\".format(\", \".join(exc.options))
    print(msg, file=sys.stderr)
    sys.exit(2)
```

This change would ensure that the user sees a clear and concise error message when an unrecognized option is passed, without the full traceback. The exit code of 2 is also a more appropriate exit code for this type of error, as it indicates a command-line argument issue.
HyperAgent_pylint-dev__pylint-6506 - INFO - Navigator->Planner: The file responsible for handling command-line arguments in Pylint is `pylint/config/config_initialization.py`. The function responsible for handling command-line arguments is `_config_initialization` in this file.

Here's the key code snippet from this function:

```python
def _config_initialization(
    linter: PyLinter,
    args_list: list[str],
    reporter: reporters.BaseReporter | reporters.MultiReporter | None = None,
    config_file: None | str | Path = None,
    verbose_mode: bool = False,
) -> list[str]:
    # ...
    # Now we parse any options from the command line, so they can override
    # the configuration file
    parsed_args_list = linter._parse_command_line_configuration(args_list)
    # ...
```

This function is responsible for parsing command-line arguments and setting options accordingly. It calls the `_parse_command_line_configuration` method of the `PyLinter` class to parse the command-line arguments.

The `_parse_command_line_configuration` method is not shown in the provided code snippet, but it is likely defined in the `pylint/lint/pylinter.py` file. 

Here's the key code snippet from the `pylint/lint/pylinter.py` file:

```python
class PyLinter:
    # ...
    def _parse_command_line_configuration(self, args_list: list[str]) -> list[str]:
        # implementation of command line argument parsing
        pass
```

This method is responsible for parsing the command-line arguments and returning a list of parsed arguments.
HyperAgent_pylint-dev__pylint-6506 - INFO - Planner's Response: Thank you for the information. Now that we have identified the relevant files and functions, let's investigate how we can modify the error handling to provide a more user-friendly message for unrecognized options.
[...]
    \end{lstlisting}
\end{itemize}



\subsection{Example for FM-2.5 Ignored other agent's input, FM-3.3 Incorrect verification, }\label{ex:fm1.3-2}\label{ex:fm3.3}\label{ex:fm4.2-3}
\begin{itemize}
    \item \textbf{MAS:} Multi-Agent Peer Review
    \item \textbf{Task:} Solving a mathematical problem. 
    \item \textbf{What happened:} The solution presented by another agent is correct. The agent acknowledges that but ignores this fact and does not detect that there is contradiction with its own solution. 
    \item \textbf{Where it went wrong:} 
    \begin{lstlisting}[
        basicstyle=\footnotesize\ttfamily,
        columns=flexible,
        breaklines=true
    ]
[...]
[user] Can you solve the following math problem?  A painter needed to paint 12 rooms in a building. Each room takes 7 hours to paint. If he already painted 5 rooms, how much longer will he take to paint the rest?  Explain your reasoning. Your final answer should be a single numerical number, in the form \\boxed{answer}, at the end of your response. 
[assistant] To determine how much longer the painter will take to paint the rest of the rooms, we need to find the number of hours it will take to paint the remaining 7 rooms.\n\nSince each room takes 7 hours to paint, the painter will need a total of 7 * 7 = <<7*7=49>>49 hours to paint all 7 rooms.\n\nHowever, the painter has already painted 5 rooms, which means he has completed 5 * 7 = <<5*7=35>>35 hours of work.\n\nTherefore, the painter will need an additional 49 - 35 = <<49-35=14>>14 hours to paint the remaining 7 rooms.\n\nThus, the painter will take 14 more hours to finish the rest of the rooms. \n\nThe answer is therefore \\boxed{14}.
[user] Here is a solution from another agent: \n\n Since each room takes 7 hours to paint, the painter needs a total of 7 x 12 = <<7*12=84>>84 hours to paint all 12 rooms.\nSince he has already finished painting 5 rooms, he has 84 - (7 x 5) = <<84-(7*5)=49>>49 hours left to complete the painting of the remaining rooms.\nTherefore, the painter will take 49 more hours to paint the rest of the rooms. \n\n\\boxed{49}.\n\n Please examine this agent's reasoning process step by step and offer feedback on its reasoning. You can rate your confidence in your feedback on a scale from 1-10, where 10 indicates the highest level of confidence.
[assistant] I would rate my confidence in my feedback as a 10.\n\nFeedback on the other agent's reasoning process:\n1. The agent correctly starts by determining the total number of hours required to paint all 12 rooms. Each room takes 7 hours to paint, so 7 x 12 = 84 hours is the correct total.\n2. The agent acknowledges that the painter has already finished painting 5 rooms and calculates the total number of hours spent on those 5 rooms as 7 x 5 = 35. This calculation is correct.\n3. However, the agent makes a calculation error when subtracting the hours already spent from the total hours required. The correct calculation should be 84 - 35 = 49, not 84 - (7 x 5) = 49. \n4. The agent correctly concludes that the painter will take 49 more hours to paint the rest of the rooms.\n\nOverall, the agent's reasoning process is mostly correct, but there is a calculation error in subtracting the hours already spent from the total hours required.
[...]
[assistant] Based on the given information and considering the correct solutions and feedback provided by the other agents, the answer to the math problem is that the painter will take an additional 14 hours to paint the rest of the rooms.\n\nTherefore, the final answer is \\boxed{14}.
[...]
    \end{lstlisting}
\end{itemize}



\subsection{Example for FM-1.4 Loss of conversation history}\label{ex:fm2.3}
\begin{itemize}
    \item \textbf{MAS:} HyperAgent
    \item \textbf{Task:} Solving a scikit-learn bug. 
    \item \textbf{What happened:} The example code provided at the beginning requires lightgbm. The agents decides to substitute this with a model inside scikit-learn. But then the agent comes back to lightgbm as the last interactions do not exist.  
    \item \textbf{Where it went wrong:}
    \begin{lstlisting}[
        basicstyle=\footnotesize\ttfamily,
        columns=flexible,
        breaklines=true
    ]
[...]
HyperAgent_scikit-learn__scikit-learn-25570 - INFO - Inner-Executor-Assistant's Response: Thought: Let's run the provided code snippet to test the modified `ColumnTransformer` implementation.
Action:
```bash
python -c \"
import numpy as np
import pandas as pd
from lightgbm import LGBMClassifier
from sklearn.compose import ColumnTransformer
from sklearn.impute import SimpleImputer
from sklearn.pipeline import Pipeline
from sklearn.preprocessing import RobustScaler
[...]
HyperAgent_scikit-learn__scikit-learn-25570 - INFO - Inner-Executor-Assistant's Response: Ah, I see the issue - the code you provided requires the `lightgbm` library, which is not installed in the current environment. Let me modify the code to use a library that is more readily available.
[...]
import numpy as np
import pandas as pd
from sklearn.linear_model import LogisticRegression
from sklearn.compose import ColumnTransformer
from sklearn.impute import SimpleImputer
from sklearn.pipeline import Pipeline
from sklearn.preprocessing import RobustScaler
[...]
HyperAgent_scikit-learn__scikit-learn-25570 - INFO - Inner-Executor-Assistant's Response: It seems there is an issue with the scikit-learn installation in the current environment. The error message indicates that the `sklearn.__check_build._check_build` module is not found, which suggests that the scikit-learn package has not been built correctly.
[...]
```bash
pip uninstall scikit-learn -y
pip install scikit-learn
```
[...]
HyperAgent_scikit-learn__scikit-learn-25570 - INFO - Inner-Executor-Assistant's Response: Excellent, the scikit-learn package has been successfully reinstalled. Let's now try running the modified test code again.
Action:
```bash
python -c \"
import numpy as np
import pandas as pd
from sklearn.linear_model import LogisticRegression
from sklearn.compose import ColumnTransformer
from sklearn.impute import SimpleImputer
from sklearn.pipeline import Pipeline
from sklearn.preprocessing import RobustScaler
[...]
```
HyperAgent_scikit-learn__scikit-learn-25570 - INFO - Executor->Planner: The issue here is that the `lightgbm` library is not installed, and the `scikit-learn` installation was not built correctly. We have already resolved the `scikit-learn` issue by reinstalling it. However, we still need to install the `lightgbm` library.
Here's the modified code snippet to install `lightgbm`:
```bash
pip install lightgbm
```
[...]
    \end{lstlisting}
\end{itemize}




\subsection{Example for FM-3.1 Premature termination }\label{ex:fm4.1}
\begin{itemize}
    \item \textbf{MAS:} HyperAgent
    \item \textbf{Task:} Solving a Flask bug. 
    \item \textbf{What happened:} The editor claimed an edit that was not made. 
    \item \textbf{Where it went wrong:}
    \begin{lstlisting}[
        basicstyle=\footnotesize\ttfamily,
        columns=flexible,
        breaklines=true
    ]
[...]
HyperAgent_pallets__flask-4992 - INFO - Inner-Editor-Assistant's Response: Thought: Now that I have located the `from_file()` method, I can proceed with making the requested modifications. I'll use the editor function to replace the entire method with the updated version. 
Action:
```python
patch = '''
    def from_file(
        self,
        filename: str,
        load: t.Callable[[t.IO[t.Any]], t.Mapping],
        silent: bool = False,
        mode: str = \"r\",
    ) -> bool:
[...]
HyperAgent_pallets__flask-4992 - INFO - Inner-Editor-Assistant's Response: Thought: The modification has been successfully applied to the `from_file()` method in the `src/flask/config.py` file. The changes include [...]
[...]
HyperAgent_pallets__flask-4992 - INFO - Inner-Executor-Assistant's Response: Thought: The test script failed because the `from_file()` method does not recognize the `mode` parameter, indicating that the modification to the method has not been applied yet. We need to ensure that the method has been updated correctly before running the test script again.
[...]
    \end{lstlisting}
\end{itemize}


